% F12：本书原创矢量示意；依据所在节的定义、证明或明确模型。
\begin{figure}[htbp]
\centering
\begin{tikzpicture}[x=1cm,y=1cm,>=stealth,every node/.style={font=\small},line width=.65pt]

\draw[->,gray](0,0)--(5,0)node[right]{\(t\)};
\draw[thick,AnalysisBlue](.4,0)--(1.5,2.4)--(2.6,0)--(3.7,2.4)--(4.8,0);
\draw[dashed,orange!80!black](.2,1.2)--(4.9,1.2);
\foreach\x in {.95,2.05,3.15,4.25}{\fill[orange!80!black](\x,1.2)circle(2pt);}
\node at(2.5,3.2){两次往返：\(N=2\)};
\draw[AnalysisBlue,very thick](8,0)--(8,2.4);
\draw[->](6,1.2)--(7.6,1.2);\node[right]at(8,1.2){同一像点的四个原像};
\node at(9.8,3.2){像集仍为 \([0,1]\)};
\node at(8.5,-.65){\(\ell(\gamma_2)=4,\quad\mathcal H^1(\gamma_2([0,1]))=1\)};

\end{tikzpicture}
\caption[参数重复经过与重数计数]{参数重复经过与重数计数。取既有往返曲线习题中的 \(N=2\)。几乎每个像点经过四次，故长度等于像上重数的积分，而非像集长度。}
\label{b1:fig:visual-f12}
\end{figure}
